\section*{Abstract}
Recent advances in long-read transcriptome sequencing enable isoform-resolved profiling in bulk, single-cell and single-nucleus samples, but partial transcripts and ambiguous read-to-isoform assignments continue to limit accurate isoform discovery and quantification, particularly when transcript annotations are incomplete. These challenges are amplified in single-cell and single-nucleus data, where coverage is sparse and biologically relevant isoforms may be restricted to specific cell populations.

We present the Long Read Alignment Assembler (LRAA), a unified framework that integrates splice-graph-based structural modeling with expectation-maximization inference to couple isoform discovery and abundance estimation. LRAA supports quantification-only, reference-guided and annotation-free modes across bulk, single-cell and single-nucleus datasets.

We benchmarked LRAA using simulated and genuine PacBio and Oxford Nanopore long-read datasets spanning sequencing standards and whole transcriptomes. We further introduced an innovative benchmark based on the Multiplexed Overexpression of Regulatory Factors (MORF) library, which provides 3k biologically expressed, barcoded isoforms with unambiguous read-level ground truth, including for partial reads. Across these evaluations, LRAA consistently outperformed most existing methods in isoform identification and quantification.

A cluster-guided single-cell workflow recovered isoforms that were low abundance in pseudobulk but enriched within specific immune and brain cell populations. In peripheral blood mononuclear cells, LRAA revealed cell-type-specific isoform usage. In frontal cortex tissue from an individual with frontotemporal dementia, LRAA detected the pathogenic cryptic spliced STMN2 isoform and differences in its prevalence across neuronal populations. Transcript-derived open reading frame predictions suggested that approximately 13\% of predicted PBMC proteoforms and 5\% of those from frontal cortex arose exclusively from unannotated splice junctions or intron retention.
